% Generated draft: requires mathematical and source review.
\documentclass[11pt]{article}
\usepackage[margin=25mm]{geometry}
\usepackage{fontspec}
\setmainfont{Latin Modern Roman}
\usepackage{amsmath,amssymb,mathtools,unicode-math}
\setmathfont{Latin Modern Math}
\usepackage{enumitem,needspace,xurl,hyperref}
\setlength{\parindent}{0pt}
\setlength{\parskip}{5pt}
\setlength{\emergencystretch}{3em}
\newcommand{\N}{\mathbb{N}}
\newcommand{\Z}{\mathbb{Z}}
\newcommand{\Q}{\mathbb{Q}}
\newcommand{\R}{\mathbb{R}}
\newcommand{\C}{\mathbb{C}}
\newcommand{\F}{\mathbb{F}}
\newcommand{\A}{\mathbb{A}}
\newcommand{\PP}{\mathbb{P}}
\newcommand{\E}{\mathbb{E}}
\newcommand{\eps}{\varepsilon}
\newcommand{\dd}{\,\mathrm{d}}
\begin{document}
\section*{UnsolvedMath ID: 701712}
\textbf{Source problem code:} catalogue entry
\subsection{Definitions and mathematical statement}
Let \(\Gamma=(V,E)\) be a finite, connected, simple, undirected graph, with graph distance \(d\) and diameter \(3\). It is distance-regular if, whenever \(d(x,y)=i\), the numbers
\[
 c_i=\bigl|\{z\in V:z\sim y,\ d(x,z)=i-1\}\bigr|,\qquad
 b_i=\bigl|\{z\in V:z\sim y,\ d(x,z)=i+1\}\bigr|
\]
depend only on \(i\), wherever the indicated distance layers exist. Put \(c_0=b_3=0\), let \(k=b_0\) be the valency, and set \(a_i=k-b_i-c_i\). The intersection array is
\[
\{b_0,b_1,b_2;c_1,c_2,c_3\}.
\]
Write the distinct eigenvalues of the adjacency matrix in decreasing order as \(k=\theta_0>\theta_1>\theta_2>\theta_3\). A distance-regular graph of diameter \(3\) is a \emph{Shilla graph} if \(\theta_1=a_3\). The supplied source states that, for a Shilla graph, \(a=a_3\) divides \(k\); it defines \(b(\Gamma)=k/a\).

The source attributes to Koolen and Park the result that there are twelve feasible intersection arrays for Shilla graphs with \(b(\Gamma)=3\). It further reports that any such graph must have one of the arrays
\[
\{12,10,3;1,3,8\},\qquad
\{12,10,5;1,1,8\},\qquad
\{15,12,6;1,2,10\}.
\]
It identifies the first two with the Doro graph and the nonisotropic unitary graph for \(q=4\), respectively. These are supplied background assertions, not independently verified here; feasibility alone is not an existence assertion.

\textbf{Existence question (catalogue entry).} Does there exist a distance-regular graph \(\Gamma\) of diameter \(3\) with intersection array
\[
\{15,12,6;1,2,10\}?
\]
This is the remaining existence question in the source's discussion of Shilla graphs with \(b(\Gamma)=3\). The source also says that A. Makhnev and N. Zyulyarkina found the automorphisms of the last graph. Because existence is explicitly posed as a question, that wording is retained as a qualified source report and is not treated as establishing existence.
\subsection{Short English statement}
Determine whether a distance-regular graph of diameter \(3\) exists with intersection array \(\{15,12,6;1,2,10\}\). The supplied record places this question in the classification of Shilla graphs with \(b(\Gamma)=3\) and labels it open. That status and the background literature claims have not been independently verified.
\subsection{Sources}
\begin{itemize}
\item Local UnsolvedMath registry: permanent ID \(701712\), source problem code \texttt{catalogue entry}, title \emph{Kourovka Notebook Problem 20.54}. The supplied record labels the problem open and gives its update timestamp as \texttt{2026-10-16T14:00:00Z}; this is registry metadata, not independent literature verification.
\item \emph{The Kourovka Notebook (Unsolved Problems in Group Theory)}, problem 20.54, 20th issue (2022), also cited in the supplied record through the 21st issue edition (2026), arXiv identifier \texttt{1401.0300v48}. The supplied collection description names E. I. Khukhro and V. D. Mazurov as editors. Proposed by A. A. Makhnev and D. O. Revin in 2022.
\item Koolen and Park: credited in the supplied statement with proving that twelve feasible intersection arrays occur as candidates for Shilla graphs with \(b=3\). No paper title, publication details, or link is supplied for this attribution.
\item A. Makhnev and N. Zyulyarkina, \emph{Doklady Maths.}, volume 84, no. 1 (2011), pp. 510--514: cited by the supplied statement concerning automorphisms of the graph with intersection array \(\{15,12,6;1,2,10\}\). No article title or link is supplied, and the precise scope of the cited result has not been verified.
\end{itemize}
\end{document}
